\section{Conclusion}

This report now supports a clearer cross-task pattern: representation quality and training scale both matter, but their effects depend on task structure. When pretrained contextual models are directly compared on aligned evaluation setups, they remain the strongest finished systems. At the same time, the final Q4 and Q5 full-data runs show that supervised recurrent models become substantially more credible once they are trained on the complete exported splits rather than on tiny capped subsets.

For Question 1, DistilBERT is the strongest sentiment classifier in the final full-data comparison, reaching test accuracy and macro-F1 of 0.9122. The two TF-IDF baselines remain strong at roughly 0.88 macro-F1, while the supervised-only BiLSTM improves with more data but still trails clearly. This again shows that the main gain comes from contextual pretraining rather than from using a neural architecture in the abstract. For Question 2, BERT dominates the CoNLL-2003 comparison and achieves the strongest overall and per-entity F1, while the feature CRF remains a strong non-transformer baseline and the BiLSTM-CRF underperforms both alternatives. For Question 3, the pretrained DistilBART summarizer outperforms TextRank on every reported metric in the controlled evaluation slice and is the strongest finished summarization system in the report. For Question 4, the full-data seq2seq-plus-attention model now reaches test BLEU 0.2725, ChrF 0.5214, METEOR 0.5431, and BERTScore 0.8291 on Multi30k, which is strong enough to demonstrate usable short-caption translation and meaningful learned alignments. The exported attention heatmap and exact-match examples confirm that the model has learned real source-target structure, even though hard cases still fail sharply.

For Question 5, the final full-data LSTM run reaches validation perplexity 162.79 and test perplexity 151.87, with its best checkpoint found at epoch 14. The training curve is smooth and the generated samples show clear gains in local fluency, even though semantic drift remains visible. Across Questions 1 through 5, the most defensible project-level claim is therefore that pretrained models provide the strongest empirical results where direct comparisons exist, but that full-data supervised neural runs can still become informative and competitive baselines once the training budget is no longer artificially constrained.

The most meaningful next step is to add matched full-data pretrained references for Questions 4 and 5 so that the stronger recurrent baselines reported here can be evaluated against transformer systems under directly aligned conditions. Even without those follow-up runs, the current written sections already support the central comparative lesson of the project: the biggest gains come from representation quality and pretraining, but scaling the data budget also changes what a classical or supervised neural baseline is capable of learning.